\newpage

\begin{thebibliography}{99}
	\makeatletter
	\renewcommand{\@biblabel}[1]{}
	\makeatother
	\bibitem{}
	\noindent Ci si limita a citare i testi --- quasi tutti in lingua italiana ---
	la cui impostazione metodologica \`{e} affine a quella qui adottata. 
	Seguendo l'articolazione del volume, si elencano dapprima i
	testi sugli \emph{argomenti di base} (algebra lineare, geometria e
	analisi) richiamati nei Preliminari; quindi i testi di \emph{meccanica
		generale} e di \emph{statica dei sistemi rigidi}, cui afferiscono la
	cinematica e la statica dei corpi rigidi, le travature reticolari e i
	teoremi dei lavori virtuali; infine alcuni testi universitari di
	\emph{Scienza delle Costruzioni} e qualche manuale tecnico per la
	consultazione. All'interno di ciascun gruppo i testi sono ordinati per
	affinit\`{a} con l'impostazione qui seguita. I riferimenti relativi alla
	meccanica del solido deformabile e alla teoria della trave sono raccolti
	nel secondo volume.
	\newline
	
	\makeatletter
	\renewcommand{\@biblabel}[1]{[#1]}
	\setcounter{enumiv}{0} \makeatother
	
	
	%=========================================================
	\subsubsection{\emph{Argomenti di base}}
	%=========================================================
	\noindent Le tematiche di base di Algebra lineare e Geometria
	necessarie alla comprensione del volume sono, ad esempio, sviluppate in:
	
	\bibitem{AbateDeFabritiis}
	M. Abate, C. de Fabritiis \emph{Geometria analitica con elementi di
		algebra lineare}, McGraw-Hill Education, terza edizione, 2015.
	
	\bibitem{Vaccaro}
	G. Vaccaro, A. Carfagna, L. Piccolella {\em Lezioni di geometria e algebra lineare}, Zanichelli, 1995.
	
	\bibitem{Abeasis}
	S. Abeasis {\em Elementi di algebra lineare e geometria}, Zanichelli, 1993.
	
	\bibitem{Dicuonzo}
	V. Dicuonzo {\em Introduzione all'algebra lineare e alla geometria analitica}, Veschi, 1991.
	
	\bibitem{Gasparini}
	I. C. Gasparini {\em Strutture algebriche e operatori lineari}, Eredi Virgilio Veschi, 1975.
	
	\vspace{0.3cm}
	\noindent mentre per l'Analisi matematica si vedano, ad esempio:
	
	\bibitem{Pagani}
	C.D. Pagani, S. Salsa \emph{Analisi matematica} vol. 1 e 2, Masson, 1998.
	
	\bibitem{Ghizzetti1}
	A. Ghizzetti, F. Rosati {\em Analisi matematica}, vol. I e II,
	MASSON editoriale Veschi, 1992.
	
	\bibitem{bramanti}
	M. Bramanti, C.D. Pagani, S. Salsa \emph{Matematica, calcolo
		infinitesimale e algebra lineare}, Zanichelli, 2000.
	
	
	%=========================================================
	\subsubsection{\emph{Meccanica generale}}
	%=========================================================
	\noindent Nozioni di Meccanica generale sono trattate in:
	
	\bibitem{LCA}
	E.N.M. Cirillo \emph{Appunti delle lezioni di meccanica
		razionale per ingegneria}, Edizioni CompoMat, 2018.
	
	\bibitem{Bordoni}
	P. G. Bordoni \emph{Lezioni di meccanica razionale}, Zanichelli, 1992.
	
	\bibitem{Benvenuti}
	P. Benvenuti, G. Maschio \emph{Appunti delle lezioni di meccanica
		razionale}, Kappa Editore, Roma, 2000.
	
	\bibitem{DenHartog2}
	J. P. Den Hartog \emph{Mechanics}, Dover Publications, Inc., New York, 1961.
	
	
	%=========================================================
	\subsubsection{\emph{Statica e meccanica dei sistemi rigidi}}
	%=========================================================
	\noindent Nozioni di Meccanica orientate alla Scienza delle Costruzioni,
	con particolare riguardo alla cinematica e alla statica dei sistemi
	rigidi, sono contenute in:
	
	\bibitem{CeradiniSCR}
	G. Ceradini {\em Cinematica e statica dei sistemi rigidi}, in Scienza
	delle Costruzioni a cura di G. Ceradini, Vol. 1, ESA Roma, 1991.
	
	\bibitem{LP}
	A. Luongo, A. Paolone {\em Meccanica delle strutture: sistemi
		rigidi ad elasticit\`{a} concentrata}, CEA, seconda edizione, 2004.
	
	\bibitem{LPDN}
	A. Luongo, A. Paolone, S. Di~Nino \textit{Rigid Structures with Point-Flexibility: Elasto-statics, Dynamics, Stability and Plastic Collapse}, Springer, 2025.
	
	\bibitem{Podio2}
	P. Podio Guidugli {\em Lezioni di statica}, Aracne editrice, 2014.
	
	\bibitem{DB}
	D. Bernardini \emph{Statica - Un'introduzione alla meccanica delle strutture}, Citt\`{a}Studi Edizioni, 2009.
	
	\bibitem{Corradi1A}
	C. Comi, L. Corradi Dell'Acqua {\em Introduzione alla meccanica
		strutturale}, McGraw-Hill, 2003.
	
	\bibitem{GBGN}
	E. Guagenti, F. Buccino, E. Garavaglia, G. Novati \emph{Statica,
		fondamenti di meccanica strutturale}, McGraw-Hill, seconda edizione, 2005.
	
	
	
	%=========================================================
	\subsubsection{\emph{Testi universitari di Scienza delle Costruzioni}}
	%=========================================================
	\noindent Per un inquadramento generale della disciplina si vedano i
	seguenti testi universitari di Scienza delle Costruzioni:
	
	\bibitem{Capurso}
	M. Capurso {\em Lezioni di Scienza delle Costruzioni}, Pitagora Editrice Bologna, 1971.
	
	\bibitem{Podio1}
	P. Podio Guidugli {\em Lezioni di Scienza delle Costruzioni}, Aracne editrice, 2009.
	
	\bibitem{Benvenuto}
	E. Benvenuto {\em La Scienza delle Costruzioni e il suo sviluppo
		storico}, Manuali Sansoni Firenze, 1981.
	
	
	%=========================================================
	\subsubsection{\emph{Manuali tecnici}}
	%=========================================================
	\noindent Pu\`{o} infine essere utile cominciare a consultare qualche
	manuale tecnico quale:
	
	\bibitem{Cremonese}
	{\em Manuale di Ingegneria civile - sezione seconda}, Edizioni scientifiche A. Cremonese - Roma, 1982.
	
	

	
	
	%=========================================================
	% RIFERIMENTI DA SPOSTARE NEL VOL. 2
	% (solido deformabile, De Saint Venant, elementi finiti, plasticit\`a,
	%  teoria della trave). Lasciati qui commentati per non perderli.
	%=========================================================
	%
%\bibitem{Borisenko}
%A.I. Borisenko, I.E. Tarapov {\em Vector and tensor analysis with
%	applications}, Dover Publications, Inc., New York 1979.
%
%\bibitem{Ceradini3}
%G. Ceradini {\em La teoria della trave}, in Scienza
%delle Costruzioni a cura di G. Ceradini, Vol. 3, ESA Roma, 1991.
%
%\bibitem{Ceradini4}
%G. Ceradini {\em Teoria lineare dei sistemi di travi}, in Scienza
%delle Costruzioni a cura di G. Ceradini, Vol. 4, ESA Roma, 1991.
%
%\bibitem{Corradi1}
%L. Corradi Dell'Acqua {\em Meccanica delle strutture: 2
%	Le teorie strutturali ed il metodo degli elementi finiti}, McGraw-Hill, 1992.
%
%\bibitem{Corradi2}
%L. Corradi Dell'Acqua {\em Meccanica delle strutture: 3 La
%	valutazione della capacit\`{a} portante}, McGraw-Hill, 1994.
%
%\bibitem{Maier}
%G. Maier {\em Complementi di scienza delle costruzioni}, clup, 1976.
%
%\bibitem{LP1}
%A. Luongo, A. Paolone {\em Scienza delle costruzioni, 1. Il
%	continuo di Cauchy}, CEA, 2004.
%
%\bibitem{LP2}
%A. Luongo, A. Paolone {\em Scienza delle costruzioni, 2. Il
%	problema di De Saint Venant}, CEA, 2005.
%
%\bibitem{PVV}
%A. Paolone, F. Vestroni, S. Vidoli {\em Scienza delle costruzioni,
%	l'analisi della tensione nelle travi - un software applicativo}, CEA, 2009.
	
\end{thebibliography}